\section{Example: converting a QFT-based adder circuit to a ripple-carry adder circuit}
%
\label{App:Adders}

The following is a four-qubit example transpilation between two forms of adder circuits, the QFT-based adder circuit proposed by Draper \cite{Draper2000}, and the ripple-carry adder proposed by Takahashi et al. \cite{Takahashi2009} for performing addition modulo $2^n$. 
These circuits have the action $\ket{a}\ket{b} \mapsto \ket{a}\ket{a+b}$. 
The bit ordering of $b = b_1b_0$ has been reversed for the sake of visual clarity in the circuit diagram.

Before proceeding with the transpilation, we present some intermediate derivations.
First, we note from~\ref{App:involution_proofs} that $H^2 = Z^2 = X^2 = \mathbf{1}$ as usual. 
\begin{align}\label{eqn:gen_involutions}
    \scalebox{0.65}{\circimg{tikz-cache/adders__derivation0__Hsquared.pdf}}
    &\approx
    \scalebox{0.65}{\circimg{tikz-cache/adders__derivation0__Zsquared.pdf}}
    \approx
    \scalebox{0.65}{\circimg{tikz-cache/adders__derivation0__Xsquared.pdf}}
    \approx
    \scalebox{0.65}{\circimg{tikz-cache/adders__derivation0__identity.pdf}}
\end{align}
Second, since the QFT-based adder circuit involves some $\lambda$ generators, we note that the position of the exponent is irrelevant. 
In particular, whether the exponent is ``attached'' to the entire $\lambda$ generator or ``attached'' to just one of the constituent gates is irrelevant, as both result in the same gate.
\begin{align*}
    \scalebox{0.65}{\circimg{tikz-cache/adders__derivation1__lambda_gen1.pdf}}
    &:=
    \scalebox{0.65}{\circimg{tikz-cache/adders__derivation1__lambda_gen2.pdf}}
    \overset{\HqcCtrlMerge}{\approx}
    \scalebox{0.65}{\circimg{tikz-cache/adders__derivation1__lambda_gen3.pdf}}
    =:
    \scalebox{0.65}{\circimg{tikz-cache/adders__derivation1__lambda_gen4.pdf}}
\end{align*}
We choose the path of least notation and place all exponents at the top right of their respective $\lambda$ generators, without drawing any boxes.

Third, pushing a Hadamard gate past a $\bullet$ or $\oplus$ gate flips one symbol to the other.
\begin{align}\label{eqn:type_flip}
    \scalebox{0.65}{\circimg{tikz-cache/adders__derivation2__type_flip1.pdf}}
    &\overset{\HqcTargDef}{\approx}
    \scalebox{0.65}{\circimg{tikz-cache/adders__derivation2__type_flip2.pdf}}
    \overset{\HqcPowDag}{\approx}
    \scalebox{0.65}{\circimg{tikz-cache/adders__derivation2__type_flip3.pdf}}
\end{align}
\begin{align}\label{eqn:HqcTargDef'}
    \scalebox{0.65}{\circimg{tikz-cache/adders__derivation2__type_flip4.pdf}}
    &\overset{\ref{eqn:type_flip}}{\approx}
    \scalebox{0.65}{\circimg{tikz-cache/adders__derivation2__type_flip5.pdf}}
    \overset{\ref{eqn:gen_involutions}}{\approx}
    \scalebox{0.65}{\circimg{tikz-cache/adders__derivation2__type_flip6.pdf}}
\end{align}
Fourth, by the same argument, pushing a $\oplus$ gate past a $\bullet$ or $\circ$ gate flips one symbol to the other.
\begin{align}\label{eqn:flip}
    \scalebox{0.65}{\circimg{tikz-cache/adders__derivation3__flip_1.pdf}}
    &\overset{\HqcOPhaseDef}{\approx}
    \scalebox{0.65}{\circimg{tikz-cache/adders__derivation3__flip_2.pdf}}
    \overset{\HqcPowDag}{\approx}
    \scalebox{0.65}{\circimg{tikz-cache/adders__derivation3__flip_3.pdf}}
\end{align}
\begin{align}\label{eqn:HqcOPhaseDef'}
    \scalebox{0.65}{\circimg{tikz-cache/adders__derivation3__flip_4.pdf}}
    &\overset{\ref{eqn:flip}}{\approx}
    \scalebox{0.65}{\circimg{tikz-cache/adders__derivation3__flip_5.pdf}}
    \overset{\ref{eqn:gen_involutions}}{\approx}
    \scalebox{0.65}{\circimg{tikz-cache/adders__derivation3__flip_6.pdf}}
\end{align}
Fifth, CNOT is self-inverse as expected.
\begin{align}\label{eqn:CNOTId}
    \scalebox{0.65}{\circimg{tikz-cache/adders__derivation4__cnot_involution_1.pdf}}
    &\overset{\HqcCtrlComp}{\approx}
    \scalebox{0.65}{\circimg{tikz-cache/adders__derivation4__cnot_involution_2.pdf}}
    \overunderset{\HQCCtrlCong}{\ref{eqn:gen_involutions}}{\approx}
    \scalebox{0.65}{\circimg{tikz-cache/adders__derivation4__cnot_involution_3.pdf}}
    \overset{\HqcCtrlId}{\approx}
    \scalebox{0.65}{\circimg{tikz-cache/adders__derivation4__cnot_involution_4.pdf}}
\end{align}
Sixth, relation \HqcCtrlComp holds in any basis. In particular, it holds in the $X$ basis.
\begin{align}\label{eqn:HqcCtrlComp'}
    \scalebox{0.65}{\circimg{tikz-cache/adders__derivation5__X_ctrl_functor_1.pdf}}
    &\overset{\HqcTargDef}{\approx}
    \scalebox{0.65}{\circimg{tikz-cache/adders__derivation5__X_ctrl_functor_2.pdf}}
    \overset{\HqcCtrlComp}{\approx}
    \scalebox{0.65}{\circimg{tikz-cache/adders__derivation5__X_ctrl_functor_3.pdf}}
    \\
    &\overunderset{\HQCCtrlCong}{\ref{eqn:HqcTargDef'}}{\approx}
    \scalebox{0.65}{\circimg{tikz-cache/adders__derivation5__X_ctrl_functor_4.pdf}}
    \overunderset{\HQCCtrlCong}{\ref{eqn:HqcTargDef'}}{\approx}
    \scalebox{0.65}{\circimg{tikz-cache/adders__derivation5__X_ctrl_functor_5.pdf}}
    \overset{\ref{eqn:gen_involutions}}{\approx}
    \scalebox{0.65}{\circimg{tikz-cache/adders__derivation5__X_ctrl_functor_6.pdf}}
\end{align}
Seventh and last, relation \HqcCtrlPair holds in any basis. In particular, it holds in the $X$ basis.
\begin{align}\label{eqn:HqcCtrlPair'}
    \scalebox{0.65}{\circimg{tikz-cache/adders__derivation6__X_ctrl_complete_basis1.pdf}}
    &\overset{\ref{eqn:gen_involutions}}{\approx}
    \scalebox{0.65}{\circimg{tikz-cache/adders__derivation6__X_ctrl_complete_basis2.pdf}}
    \overset{\HqcCtrlDag}{\approx}
    \scalebox{0.65}{\circimg{tikz-cache/adders__derivation6__X_ctrl_complete_basis3.pdf}}
    \overset{\HqcTargDef,\HqcTargMDef}{\approx}
    \scalebox{0.65}{\circimg{tikz-cache/adders__derivation6__X_ctrl_complete_basis4.pdf}}
\end{align}

With those preparatory derivations out of the way, we proceed to the transpilation. We begin from a QFT-based adder circuit \cite{Draper2000} and show an explicit step-by-step transpilation to a ripple-carry adder circuit \cite{Takahashi2009}. All intermediate circuits are semantically equivalent by soundness of Fig. \ref{Fig:HQCEqnTheory}, meaning they all implement the same unitary transformation.
\begin{align*}
    \scalebox{0.65}{\circimg{tikz-cache/adders__adders1.pdf}}
    &:=
    \scalebox{0.65}{\circimg{tikz-cache/adders__adders2.pdf}}
    \\
    \overunderset{\HQCCtrlCong}{\ref{eqn:HqcTargDef'}}{\approx}
    \scalebox{0.65}{\circimg{tikz-cache/adders__adders3.pdf}}
    &\overunderset{\HQCCtrlCong}{\ref{eqn:HqcTargDef'}}{\approx}
    \scalebox{0.65}{\circimg{tikz-cache/adders__adders4.pdf}}
    \\
    \overunderset{\HQCCtrlCong}{\ref{eqn:HqcTargDef'}}{\approx}
    \scalebox{0.65}{\circimg{tikz-cache/adders__adders5.pdf}}
    &\overunderset{\HQCCtrlCong}{\ref{eqn:HqcTargDef'}}{\approx}
    \scalebox{0.65}{\circimg{tikz-cache/adders__adders6.pdf}}
    \\
    \overset{\ref{eqn:gen_involutions}}{\approx}
    \scalebox{0.65}{\circimg{tikz-cache/adders__adders7.pdf}}
    &\overset{\HqcCtrlComp}{\approx}
    \scalebox{0.65}{\circimg{tikz-cache/adders__adders8.pdf}}
    \approx
    \scalebox{0.65}{\circimg{tikz-cache/adders__adders9.pdf}}
    \\
    \overset{\HqcCtrlDag}{\approx}
    \scalebox{0.65}{\circimg{tikz-cache/adders__adders10.pdf}}
    &\overunderset{\HQCCtrlCong}{\ref{eqn:HqcOPhaseDef'}}{\approx}
    \scalebox{0.65}{\circimg{tikz-cache/adders__adders11.pdf}}
    \overunderset{\HQCCtrlCong}{\HqcCtrlPair}{\approx}
    \scalebox{0.65}{\circimg{tikz-cache/adders__adders12.pdf}}
    \\
    \overunderset{\HQCCtrlCong}{\HqcPowAdd}{\approx}
    \scalebox{0.65}{\circimg{tikz-cache/adders__adders12a.pdf}}
    &\overunderset{\HQCCtrlCong}{\HqcPowId,\HqcPowZero}{\approx}
    \scalebox{0.65}{\circimg{tikz-cache/adders__adders13.pdf}}
    \overset{\HqcCtrlComp}{\approx}
    \scalebox{0.65}{\circimg{tikz-cache/adders__adders14.pdf}}
    \\
    \overset{\ref{eqn:CNOTId}}{\approx}
    \scalebox{0.65}{\circimg{tikz-cache/adders__adders15.pdf}}
    &\approx
    \scalebox{0.65}{\circimg{tikz-cache/adders__adders16.pdf}}
    \overset{\ref{eqn:HqcCtrlComp'}}{\approx}
    \scalebox{0.65}{\circimg{tikz-cache/adders__adders17.pdf}}
    \\
    \overunderset{\HQCCtrlCong}{\ref{eqn:HqcCtrlPair'}}{\approx}
    \scalebox{0.65}{\circimg{tikz-cache/adders__adders18.pdf}}
    &\overunderset{\HQCCtrlCong}{\HqcTargDef}{\approx}
    \scalebox{0.65}{\circimg{tikz-cache/adders__adders18a.pdf}}
    \overunderset{\HQCCtrlCong}{\HqcCtrlDag}{\approx}
    \scalebox{0.65}{\circimg{tikz-cache/adders__adders18b.pdf}}
    \\
    \overunderset{\HQCCtrlCong}{\ref{eqn:gen_involutions}}{\approx}
    \scalebox{0.65}{\circimg{tikz-cache/adders__adders19.pdf}}
    &\overset{\HqcCtrlDag}{\approx}
    \scalebox{0.65}{\circimg{tikz-cache/adders__adders20.pdf}}
    =: 
    \begin{array}{c}
         \text{Takahashi} \\
         \text{ripple-carry} \\
         \text{adder } \cite{Takahashi2009}
    \end{array}
\end{align*}